% Section "nn" (writer: oodcomp). Source: research/tracks/context/notes-final.md §4 (Props 4.1--4.5, §4.3, Thm 4.6,
% Cor 4.7, Rem 4.8, §4.5--4.7, §7 items 6, 7), the Vaintrob--Hänni slides (research/prior/), and Hänni's notes
% hanni-kolmogorov-prior-in-NNs.md, hanni-bridging-parametrized-bayes-and-solomonoff.md. Renamings (NOTATION.md):
% kernel profile kappa -> \Psi, kernel k -> \mathcal K, eigenvalue -> \Lambda_S, parameter prior pi -> \varpi,
% ridge lambda -> \beta, circuit size c(f) -> \csize(f). The slides' kappa_N, C, C_L are quoted in their notation.
\section{Neural-network Bayes and circuit SFI}\label{sec:nn}

H\"anni describes Bayesian learning with neural networks (NN Bayes) as ``kinda like a circuit version of solomonoff function induction'' \citep{hanni2026request}, referring to slides by Vaintrob and H\"anni \citep{vaintrobhanni2026nnbayes}. This section states what is proved about this link, what the slides claim, and what is open. An \emph{architecture} is a parameter space $\Theta$ with a prior $\varpi$ and a map $\theta\mapsto f_\theta:\{0,1\}^d\to Y$.

\begin{proposition}[NN Bayes is function-space Bayes]\label{prop:nn:functionspace}
\status{proved (trivial)}\src{context Prop 4.1}
If the likelihood depends on $\theta$ only through the values of $f_\theta$ on the data (a 0--1 likelihood, a margin rule, any label-noise model), the network's posterior predictive equals the function-space Bayes predictive with prior $P_A(f):=\varpi\{\theta:f_\theta=f\}$.
\end{proposition}

\emph{Proof.} Group the parameter integral by the value of $f_\theta$. \qed

So NN Bayes is function induction with the prior $P_A$, and the question is how $P_A$ compares with $2^{-K}$, with $\MF$, or with a circuit prior.

\begin{proposition}[coding-theorem upper bound]\label{prop:nn:upper}
\status{proved (given the coding theorem)}\src{context Prop 4.2}
If $P_A$, as a semimeasure on truth tables, is lower-semicomputed by a program $a$ (encoding the architecture, the prior and $d$), then $P_A(f)\le2^{-K(f\mid a)+c}\le2^{-K(f)+K(a)+c'}$, with $c$ independent of the architecture.
\end{proposition}

This is the conditional coding theorem for the family indexed by $a$. $P_A$ is lsc for finite parameter sets with rational probabilities; for Gaussian priors on real weights with sign outputs we only sketch it (up to a null boundary, $\{\theta:f_\theta=f\}$ is an open semialgebraic set, approximable from below by certified rational boxes). H\"anni's version has an extra $\log\log(1/p)$ term \citep{hanni2026notes}, which is not needed. The bound is an exact, uncomputable form of the simplicity-bias bound of \citet{dingle2018input}.

\begin{proposition}[every prior nearly saturates the bound on its typical set]\label{prop:nn:typical}
\status{proved; known (standard)}\src{context Prop 4.3}
For every semimeasure $P$ on truth tables, every string $a$ and $j\ge0$: $P\{f:P(f)\le2^{-K(f\mid a)-j}\}\le2^{-j}$. Hence if $P_A$ is lower-semicomputed by $a$ and $\varpi$ is a probability measure, then with probability at least $1-2^{-j}$ over $\theta\sim\varpi$, $K(f_\theta\mid a)-c\le-\log P_A(f_\theta)<K(f_\theta\mid a)+j$.
\end{proposition}

\emph{Proof.} The left side is at most $\sum_f2^{-K(f\mid a)-j}\le2^{-j}$ by Kraft's inequality; the lower bound is \cref{prop:nn:upper}. \qed

This is the standard fact that $-\log P\approx K$ on $P$-most objects \citep{li2008introduction} (cf.\ \cref{lem:ex:surprisal}). It confirms H\"anni's remark that the result ``should just be true about any low complexity prior on functions --- the NNs aren't doing anything here'': $P_A$ looks like $2^{-K}$ on its own typical set, which can omit most simple functions.

\begin{proposition}[no general lower bound]\label{prop:nn:nolower}
\status{proved}\src{context Prop 4.4}
(a) Let $F_d$ be the set of functions an architecture $A_d$ can compute, with $|F_d|<2^{2^d}$ and membership decidable given $d$ (any finite parameter set; any real-weight threshold architecture, by Tarski's decision procedure). The first truth table $f_d\notin F_d$ has $K(f_d)\le K(d)+c$ and $P_{A_d}(f_d)=0$. So no bound $P_A(f)\ge2^{-aK(f)-b}$ with $a,b$ independent of $d$ holds for such a family. (b) The parameter--function map alone implies no simplicity bias: a lookup-table architecture with uniform parameters gives every $f$ probability $2^{-2^d}$.
\end{proposition}

The slides make the point of (a) with functions that diagonalise against all small circuits. So the right comparison object for networks is \emph{circuit} SFI, not program SFI. For the uniform prior over circuits of one size, one half of the comparison holds (\cref{def:comp:circuits} for $N_s$ and $\csize$):

\begin{proposition}[uniform circuit priors contain the circuit simplicity prior]\label{prop:nn:uniform}
\status{proved}\src{context Prop 4.5}
Let $P_{{\rm unif},S}$ be uniform over circuits of size exactly $S$. For every $f$ with $\csize(f)<S$, $P_{{\rm unif},S}(f)\ge1/(N_{\csize(f)}\cdot16(d+1+S)^2)$. Hence Bayes with this prior makes at most $(\ln2)(\log N_{\csize(f)}+2\log(d+1+S)+4)=O(\csize(f)\log(\csize(f)+d)+\log S)$ Gibbs mistakes on every ordering.
\end{proposition}

\emph{Proof idea} (\cref{app:nn}). Pad a minimal circuit with arbitrary dead gates and a final gate copying its output. This is one half of H\"anni's ``first main afaik-open problem'' \citep{hanni2026notes}; the other half, that the uniform prior is not much \emph{larger} than the simplicity prior on functions without small circuits, is open.

\begin{remark}[what a bridge would give]\label{rem:nn:bridge}
\status{remark}\src{context \S4.3}
If $P_A(f)\ge2^{-a\ell(f)-b}$ for a description length $\ell$, then on every ordering NN Bayes makes at most $(\ln2)(a\ell(f)+b)$ Gibbs mistakes and at most $a\ell(f)+b$ posterior-majority mistakes (as in \cref{thm:ood:cumulative,rem:comp:majority}); if also $P_A(f)\le2^{-a'\ell(f)+b'}$, the log loss on the full domain is at least $a'\ell(f)-b'$. \Cref{prop:nn:uniform} is the lower half for uniform circuit priors. The converse, ``few mistakes implies a small circuit'', which H\"anni thought ``true for NN bayes as well'', is claimed by the slides in their setting and in distribution (\cref{rem:nn:slides}), and holds for circuit Bayes on every ordering (\cref{thm:comp:circuitbayes}).
\end{remark}

\paragraph{The Gaussian-process limits.} H\"anni suggests that in the infinite-width Gaussian limits (NNGP and NTK) the network is ``doing sth more like a polyfit'' than circuit SFI \citep{hanni2026notes}. On the cube $\{\pm1\}^d$ with uniform measure, let $\chi_S(x)=\prod_{j\in S}x_j$ be the Walsh characters and, for a kernel $\mathcal K$, $(T_{\mathcal K}g)(x):=\Ex_{x'}\mathcal K(x,x')g(x')$.

\begin{theorem}[Walsh diagonalisation; dimension bound]\label{thm:nn:walsh}
\status{proved; (a) known, (b) standard; computed (context c2)}\src{context Thm 4.6}
(a) If $\mathcal K(x,x')=\Psi(\langle x,x'\rangle/d)$, every $\chi_S$ is an eigenfunction of $T_{\mathcal K}$, with eigenvalue $\Lambda_{|S|}$ depending only on $|S|$. This covers the NNGP and NTK kernels of fully connected networks with i.i.d.\ Gaussian weights and biases, which are functions of $\langle x,x'\rangle$ on the cube.
(b) If $\mathcal V$ is any space of functions on the cube of dimension at most $n$ (for instance $\mathrm{span}\{\mathcal K(x_i,\cdot)\}_{i\le n}$, which contains the GP posterior mean and the kernel ridge regression solution), then $\binom dk^{-1}\sum_{|S|=k}\min_{g\in\mathcal V}\|\chi_S-g\|^2\ge1-n/\binom dk$.
\end{theorem}

\emph{Proof idea} (\cref{app:nn}). (a) $\mathcal K(x,x')$ depends only on $x\odot x'$, so characters are eigenfunctions; permutation symmetry makes the eigenvalue depend on $|S|$ only. (b) The squared projections of all characters onto $\mathcal V$ sum to $\dim\mathcal V$. Part (a) is in \citet{yang2019fine} (as summarised in its abstract and in the slides) and is proved in the slides by the same argument; (b) is the standard dimension argument, cf.\ \citet{kamath2020approximate}.

\begin{corollary}[kernels versus SFI on sparse parities]\label{cor:nn:separation}
\status{proved; computed (context c2)}\src{context Cor 4.7}
For $|S|=k$ and i.i.d.\ uniform inputs, SFI ($\nu\in\{W,\MF\}$) has cumulative posterior error at most $(\ln2)(K(\chi_S)+c)$ on every sequence of inputs, with $K(\chi_S)\le\log\binom dk+2\log(k+1)+2\log(d+1)+c$, and the squared error of its posterior-mean predictor averaged over the first $n$ steps is at most $4(\ln2)(K(\chi_S)+c)/n$. Every kernel method or GP posterior mean needs $n\ge(1-\varepsilon)\binom dk$ samples for average squared error at most $\varepsilon$ over random $k$-sparse targets, while SFI needs $O(k\log d/\varepsilon)$ in this averaged sense.
\end{corollary}

The separation is in squared error of the real-valued predictor; we do not claim a 0--1 version in general. For one kernel family a 0--1 failure is exact:

\begin{remark}[bias-free two-layer ReLU kernels and odd parities]\label{rem:nn:oddparity}
\status{proved; computed (context c2, c10, c12); known}\src{context Rem 4.8(a)--(c)}
(a) For the bias-free two-layer ReLU NNGP, $\Psi_1(\rho)=(\sqrt{1-\rho^2}+\rho(\pi-\arccos\rho))/\pi$, and for the NTK, $\Psi_1(\rho)+\rho\Psi_0(\rho)$ with $\Psi_0(\rho)=(\pi-\arccos\rho)/\pi$, the function $\Psi(\rho)-a\rho$ is even ($a=\frac12$, resp.\ $a=1$). Hence $\Lambda_{|S|}=0$ for every odd $|S|\ge3$, $\mathcal V:=\mathrm{span}\{\mathcal K(x_i,\cdot)\}\perp\chi_S$, and $\|\chi_S-\hat g\|^2=1+\|\hat g\|^2\ge1$ for every $n$: the kernel predictor never learns these parities. This is the cube version of the finding of \citet{basri2019convergence} that a shallow network without bias cannot learn odd frequencies (on the sphere, as stated in their abstract). Biases or a third layer make the eigenvalues nonzero.
(b) Let $z\sim\mathrm{GP}(0,\mathcal K)$ for one of these kernels and $|S|\ge3$ odd. Then $\Ex_xz(x)\chi_S(x)$ is Gaussian with variance $\Lambda_{|S|}=0$, so it vanishes almost surely, while $\mathrm{sign}\,z=\chi_S$ would force it to equal $\Ex|z|>0$. So the GP prior, and every posterior obtained by conditioning on an event of positive probability (sign or margin likelihoods), gives probability exactly $0$ to the classifier $\chi_S$.
(c) For distinct training inputs $x_1,\dots,x_n$ with labels $\chi_S(x_i)$ and the set $U$ of the other $2^d-n$ points, every $g\in\mathcal V$ has $\sum_{x\in U}g(x)\chi_S(x)=-\sum_{i\le n}g(x_i)\chi_S(x_i)$. For kernel ridge regression (the GP mean with noise variance $\beta>0$) the training sum is $y^\top K(K+\beta I)^{-1}y\ge0$, so the predictor is weakly anti-correlated with the target on the unseen points.
A direct computation (\cref{app:nn}) gives kernel ridge regression a sign accuracy on unseen points of $0.05$ to $0.07$ at $n=900$, $d=10$, far below chance; this is computed, not proved.
\end{remark}

\begin{remark}[the Vaintrob--H\"anni slides]\label{rem:nn:slides}
\status{known; unverified (claimed in the slides); warm-up arithmetic re-checked (context c8)}\src{context \S4.5; slides 10--17, 29--37, App.\ A, F, J, K}
The slides warn that their appendices ``have largely been generated by Claude'' and have not been ``meticulously reviewed'', that the main results summarise ``much more caveated and detailed versions'', that only the strong-prior case is proved in the talk, and that Bayesian learning and SGD are ``importantly different''. They claim: (i) for programs with prior $\propto3^{-|p|}$ and i.i.d.\ inputs, $O(C)$ samples ($C$ the length of the shortest program for $f$) give average posterior probability above $99.9\%$ on the correct answer, with probability $1-\exp(-\Omega(C))$, and a ``relevance-based'' sample-complexity bound for any countable class; (ii) for an MLP of depth $L$ and any width $N\ge C:=C_L(f)$ (the least width and weight bound of a depth-$L$ MLP implementing $f$), with bounded, Lipschitz, non-polynomial activation and Gaussian weights in the strong-prior limit, $\mathrm{poly}(d,C)$ samples suffice in the same sense, the polynomial hiding a factor exponential in depth; (iii) the orderings of functions by $C_L$, by a layer-weighted squared norm $\kappa_N$ and by prior mass are ``scrambled only polynomially'', uniformly in width; (iv) an in-distribution mistake bound $\mathrm{poly}(d,C_L(f))$ in expectation, and ``the circuit complexity $C_L$ of the learned function is at most $\mathrm{poly}(d,\#\text{mistakes made while learning it})$''; (v) a mean-field analogue \citep{vaintrob2026width}, while the NNGP ``does not have a good circuit prior''; (vi) out of distribution the guarantee ``is simply false'' (\cref{thm:comp:ladder}). We re-checked the proof and arithmetic of (i), the character argument of their Appendix J (\cref{thm:nn:walsh}(a)) and the ladder of Appendix K; the rest is unchecked. \Cref{cor:learn:perstep}(a) is the Markov half of the relevance-based bound, applied to $W$ and $\MF$. Claim (iii) concerns orderings, not two-sided bounds on prior mass; the slides list ``a lower bound on the log prior of a function'' as open.
\end{remark}

\begin{remark}[what can be said rigorously]\label{rem:nn:rigorous}
\status{remark}\src{context \S4.6}
Proved: NN Bayes is function-space Bayes; its prior is at most $2^{-K+O(K(a))}$, nearly equal on its typical set, with no general lower bound, so the comparison object is circuit SFI; uniform circuit priors contain the circuit simplicity prior; the Gaussian-process limits fail on sparse and odd parities where SFI succeeds; and for circuit Bayes, mistakes and circuit size are polynomially related on every ordering (\cref{thm:comp:circuitbayes}), while Occam selection on the same prior fails exponentially out of distribution (\cref{thm:comp:ladder}). Claimed in the slides: the strong-prior theorems. Open: two-sided bounds on $P_A(f)$ at fixed $\sigma$, and whether NN Bayes generalises like SGD-trained networks at all, which H\"anni doubts \citep{hanni2026notes}.
\end{remark}

\begin{remark}[simplicity-bias literature]\label{rem:nn:literature}
\status{known (snippets; unverified)}\src{context \S4.7}
\citet{valle2019deep} argue, following \citet{dingle2018input}, that the parameter--function map of deep networks is exponentially biased towards simple functions. H\"anni's objection \citep{hanni2026notes} that the upper bound holds for every map while the needed lower bound fails is made rigorous by \cref{prop:nn:typical,prop:nn:nolower}; the slides argue the same. \citet{mingard2021sgd} report that SGD samples functions with probabilities close to the Bayesian posterior, and \citet{mingard2025deep} attribute much of generalisation to a prior bias towards simple functions; we take no position. \citet{buzaglo2024} show that random interpolators under a flat parameter prior generalise when a narrow teacher exists; the slides call this the closest prior work and note that its bound degrades with student width. \citet{goldblum2023no} relate no-free-lunch theorems to Kolmogorov complexity, and \citet{grau2024learning} meta-train networks on data from universal machines to approximate Solomonoff induction. These descriptions follow the abstracts.
\end{remark}

\paragraph{Open questions.}\label{q:nn:bridges} The upper half of the uniform-circuit bridge (\cref{prop:nn:uniform}); a lower bound on the NN log prior at fixed $\sigma$; a check of the slides' main theorem and of ``mistakes imply a small circuit''; the per-point test error of GP classification on odd parities; a proof that kernel-ridge sign accuracy on unseen points falls below $\frac12$.
